\section{基本算法与杂项}

蔡勒公式
\[
w=\left(y+\left[{\frac {y}{4}}\right]+\left[{\frac {c}{4}}\right]-2c+\left[{\frac {26(m+1)}{10}}\right]+d-1\right){\bmod {7}}
\]
或
\[
w=\left(y+\left[{\frac {y}{4}}\right]+\left[{\frac {c}{4}}\right]-2c+2m+\left[{\frac {3(m+1)}{5}}\right]+d+1\right){\bmod {7}}
\]

\subsection{三分}

A. 整数域三分，求凸函数的极小值
\begin{lstlisting}[language = C++]
ll ternary_search(ll l,ll r){
	ll lmid, rmid;
	while(r-l > 2){
		lmid = l+(r-l)/3;
		rmid = r-(r-l)/3;
		if(check(lmid) <= check(rmid)) r = rmid;
		else l = lmid; 
	}
	return min(check(l), min(check(l+1), check(l+2)));
} 
\end{lstlisting}

B. 实数域三分，求凹函数的极大值
\begin{lstlisting}[language = C++]
double ternary_search(double l, double r){
	double mid, lmid, rmid;
	while(r-l > eps){
        mid = l+(r-l)/2;
        lmid = mid-eps, rmid = mid+eps;
        if(f(lmid) > f(rmid)) r = mid;
        else l = mid;
    }
	return f(l);
}
\end{lstlisting}

\subsection{对拍}

\begin{lstlisting}[language=C++]
// Windows
void solve() {
    system("g++ -std=c++23 data.cpp -o data.exe");
    system("g++ -std=c++23 std.cpp -o std.exe");
    system("g++ -std=c++23 solve.cpp -o solve.exe");
    int t = 0;
    while (true) {
        cout << "test" << ++t << ": ";
        system("data.exe > data.in");
        system("std.exe < data.in > std.ans");
        system("solve.exe < data.in > solve.ans");
        if (system("fc std.ans solve.ans > diff.log")) {
            cout << "WA" << endl;
            system("type data.in");
            system("type std.ans");
            system("type solve.ans");
            system("pause");
            break;
        }
        cout << "AC" << endl;
    }
}
// Linux
void solve() {
    system("rm data; rm solve; rm std; make data && make solve && make std");
    int t = 0;
    while (true) {
        cout << "test" << ++t << ": ";
        system("./data > in");
        system("./solve < in > solve.ans");
        system("./std < in > std.ans");
        if (system("diff solve.ans std.ans")) {
            cout << "WA" << endl;
            system("pause");
            break;
        }
        cout << "AC" << endl;
    }
}
\end{lstlisting}

\subsection{高精度}

\begin{lstlisting}[language=C++]
struct BigInteger {
    vector<int> digits;
    bool isNegative = false;

    BigInteger(int num) {
        if (num < 0) {
            isNegative = true;
            num = -num;
        }
        do {
            digits.push_back(num % 10);
            num /= 10;
        } while (num > 0);
        removeLeadingZeros();
    }
    BigInteger(const string& num = "0") {
        if (num[0] == '-') {
            isNegative = true;
            for (int i = num.size() - 1; i > 0; i--) {
                digits.push_back(num[i] - '0');
            }
        } else {
            for (int i = num.size() - 1; i >= 0; i--) {
                digits.push_back(num[i] - '0');
            }
        }
        removeLeadingZeros();
    }

    BigInteger operator+(const BigInteger& other) const {
        if (isNegative == other.isNegative) {
            BigInteger result;
            result.isNegative = isNegative;
            int carry = 0;
            for (size_t i = 0; i < max(digits.size(), other.digits.size()) || carry; i++) {
                if (i < digits.size()) carry += digits[i];
                if (i < other.digits.size()) carry += other.digits[i];
                result.digits.push_back(carry % 10);
                carry /= 10;
            }
            result.removeLeadingZeros();
            return result;
        }
        return isNegative ? (other - (-*this)) : (*this - (-other));
    }
    BigInteger operator-(const BigInteger& other) const {
        if (isNegative == other.isNegative) {
            if (isNegative) return (-other) - (-*this);
            if (*this < other) return -(other - *this);
            BigInteger result;
            int borrow = 0;
            for (size_t i = 0; i < digits.size(); ++i) {
                int diff = digits[i] - borrow;
                if (i < other.digits.size()) diff -= other.digits[i];
                if (diff < 0) {
                    diff += 10;
                    borrow = 1;
                } else {
                    borrow = 0;
                }
                result.digits.push_back(diff);
            }
            result.removeLeadingZeros();
            return result;
        }
        return isNegative ? -(-*this + other) : (*this + (-other));
    }
    BigInteger operator*(const BigInteger& other) const {
        BigInteger result;
        result.digits.resize(digits.size() + other.digits.size());
        for (size_t i = 0; i < digits.size(); ++i) {
            int carry = 0;
            for (size_t j = 0; j < other.digits.size() || carry; ++j) {
                int cur = result.digits[i + j] + digits[i] * (j < other.digits.size() ? other.digits[j] : 0) + carry;
                result.digits[i + j] = cur % 10;
                carry = cur / 10;
            }
        }
        result.isNegative = (isNegative != other.isNegative);
        result.removeLeadingZeros();
        return result;
    }
    BigInteger operator/(int divisor) const {
        if (divisor == 0) {
            throw runtime_error("Division by zero");
        }

        bool resultNegative = (isNegative ^ (divisor < 0));
        BigInteger dividend = *this;
        dividend.isNegative = false;
        divisor = abs(divisor);
        if (dividend.digits[0] == 0 && dividend.digits.size() == 1) {
            return BigInteger("0");
        }
        string quotientStr;
        int remainder = 0;
        for (int i = dividend.digits.size() - 1; i >= 0; i--) {
            remainder = remainder * 10 + dividend.digits[i];
            int q = remainder / divisor;
            remainder %= divisor;
            quotientStr += to_string(q);
        }
        size_t firstNonZero = 0;
        while (firstNonZero < quotientStr.size() && quotientStr[firstNonZero] == '0') {
            firstNonZero++;
        }
        quotientStr = (firstNonZero == quotientStr.size()) ? "0" : quotientStr.substr(firstNonZero);
        if (resultNegative && remainder != 0) {
            BigInteger temp(quotientStr);
            temp = temp + BigInteger("1"); // 正数加 1，转为负数后即向下取整
            temp.isNegative = true;
            return temp;
        }
        BigInteger quotient(quotientStr);
        quotient.isNegative = resultNegative;
        return quotient;
    }
    bool operator==(const BigInteger& other) const {
        if (isNegative != other.isNegative) {
            return false;
        }
        if (digits.size() != other.digits.size()) {
            return false;
        }
        for (size_t i = 0; i < digits.size(); ++i) {
            if (digits[i] != other.digits[i]) {
                return false;
            }
        }
        return true;
    }
    bool operator>(const BigInteger& other) const {
        if (isNegative != other.isNegative) {
            return isNegative;
        }
        if (digits.size() != other.digits.size()) {
            return (digits.size() > other.digits.size()) ^ isNegative;
        }
        for (int i = digits.size() - 1; i >= 0; i--) {
            if (digits[i] != other.digits[i]) {
                return (digits[i] > other.digits[i]) ^ isNegative;
            }
        }
        return false;
    }
    bool operator<(const BigInteger& other) const {
        if (isNegative != other.isNegative) {
            return isNegative;
        }
        if (digits.size() != other.digits.size()) {
            return (digits.size() < other.digits.size()) ^ isNegative;
        }
        for (int i = digits.size() - 1; i >= 0; i--) {
            if (digits[i] != other.digits[i]) {
                return (digits[i] < other.digits[i]) ^ isNegative;
            }
        }
        return false;
    }
    BigInteger operator-() const {
        BigInteger result = *this;
        result.isNegative = !isNegative;
        return result;
    }
    void removeLeadingZeros() {
        while (digits.size() > 1 && digits.back() == 0) {
            digits.pop_back();
        }
        if (digits.size() == 1 && digits[0] == 0) {
            isNegative = false;
        }
    }
    friend istream& operator>>(istream& is, BigInteger& num) {
        string input;
        is >> input;
        num = BigInteger(input);
        return is;
    }
    friend ostream& operator<<(ostream& os, const BigInteger& num) {
        if (num.isNegative && !(num.digits.size() == 1 && num.digits[0] == 0)) {
            os << '-';
        }
        for (int i = num.digits.size() - 1; i >= 0; i--) {
            os << num.digits[i];
        }
        return os;
    }
};
\end{lstlisting}

\subsection{日期类}

\begin{lstlisting}[language=C++]
int daysInM[13] = {0, 31, 28, 31, 30, 31, 30, 31, 31, 30, 31, 30, 31};
int daysBeforeM[13] = {0, 0, 31, 59, 90, 120, 151, 181, 212, 243, 273, 304, 334};
bool isLeapYear(int year) {
    return year % 400 == 0 || year % 4 == 0 && year % 100 != 0;
}
int daysInMonth(int year, int month) {
    return daysInM[month] + (month == 2 ? isLeapYear(year) : 0);
}
struct Date {
    int year, month, day;
    bool isLeap;
    int daysFromAD; // 公元1年1月1日(星期一)到当前日期总天数，算上了1582.10.5~14
    int weekDay;
    Date(int year_, int month_, int day_) : year(year_), month(month_), day(day_) {
        isLeap = isLeapYear(year);
        int y = year - 1;
        daysFromAD = y * 365 + y / 4 - y / 100 + y / 400 + daysBeforeM[month] + day - 1;
        if (isLeap && month > 2) {
            daysFromAD++;
        }
        weekDay = (daysFromAD + 1) % 7;
    }
    Date(int days) {
        daysFromAD = days;
        weekDay = (daysFromAD + 1) % 7;
        year = days / 365 + 1;
        int d;
        while (year > 1) {
            int y = year - 1;
            d = y * 365 + y / 4 - y / 100 + y / 400;
            if (d <= days) {
                break;
            }
            year--;
        }
        isLeap = isLeapYear(year);
        d = days - d + 1;
        month = 1;
        while (month < 12) {
            int dIM = daysInMonth(year, month);
            if (d <= dIM) {
                break;
            }
            d -= dIM;
            month++;
        }
        day = d;
    }
    void print() {
        cout << year << "/" << month << "/" << day << ", " << weekDay << ", " << daysFromAD << "\n";
    }
};
int daysBetween(const Date &u, const Date &v) {
    return abs(u.daysFromAD - v.daysFromAD);
}
\end{lstlisting}

\subsection{取模类}

\begin{lstlisting}[language=C++]
using ll = long long;

template<typename T>
T fpow(T x, int y){
    T res = 1;
    for(; y; y >>= 1){
        if(y&1) res *= x;
        x *= x;
    }
    return res;
}

template<int M>
struct Mint{
    ll val;
    Mint(): val(0) {}
    Mint(ll x): val((x%M+M)%M) {}
    Mint(const Mint<M> &t): val(t.val) {}
    Mint inv() const {return fpow<Mint<M>>(*this, M-2);}
    Mint inv() {return fpow<Mint<M>>(*this, M-2);}
    Mint<M>& operator = (const ll &t){
        val = (t%M+M)%M;
        return *this;
    }
    Mint<M>& operator += (const Mint<M> &t){
        val += t.val;
        if(val >= M) val -= M;
        return *this;
    }
    Mint<M>& operator -= (const Mint<M> &t){
        val -= t.val;
        if(val < 0) val += M;
        return *this;
    }
    Mint<M>& operator *= (const Mint<M> &t){
        val = (val*t.val)%M;
        return *this;
    }
    Mint<M>& operator /= (const Mint<M> &t){
        (*this) *= t.inv();
        return *this;
    }
    Mint<M> operator + (const Mint<M> &t){
        Mint<M> u(*this);
        return u += t;
    }
    Mint<M> operator - (const Mint<M> &t){
        Mint<M> u(*this);
        return u -= t;
    }
    Mint<M> operator * (const Mint<M> &t){
        Mint<M> u(*this);
        return u *= t;
    }
    Mint<M> operator / (const Mint<M> &t){
        Mint<M> u(*this);
        return u /= t;;
    }
    bool operator != (const Mint<M> &t) {return t.val != val;}
    bool operator != (ll &t) {return t != val;}
    bool operator == (const Mint<M> &t) {return t.val == val;}
    bool operator == (ll &t) {return t == val;}
    friend ostream& operator << (ostream &os, const Mint<M> &t){
        os << t.val;
        return os;
    }
    friend istream& operator >> (istream &is, Mint<M> &t){
        ll num; is >> num;
        t = num;
        return is;
    }
};
\end{lstlisting}





